\chapter{Auswertung von Termbildmengen}
\section{Einsetzung erzeugender Terme}

Eine Zwischenmenge kann durch ihre erzeugenden Terme ausgewertet werden:
Ein Zeuge aus dieser Menge wird durch die Indizes ersetzt, die ihn erzeugen.
Der folgende Satz gilt für einen beliebigen Mengenterm \(s(i,j)\) und
ein beliebiges Prädikat \(F(u)\), jeweils mit festen Parametern.
Die verschiedenen Bindungsvariablen \(u,i,j\) kommen nicht in
\(U,I,J\) oder in diesen festen Parametern frei vor; \(u\) kommt auch
in \(s(i,j)\) nicht frei vor. Die Einsetzung \(u:=s(i,j)\) in \(F\)
ist kollisionsfrei. Insbesondere werden gebundene Variablen in \(F\)
nötigenfalls vorher umbenannt. Die Indexmengen dürfen leer sein.

\FormulaThmDeltaK[Auswertung einer durch Terme dargestellten Zwischenmenge]{%
  \begin{aligned}[t]
    &\forall u\,\bigl(u\in U\leftrightarrow
      \exists i\in I\,\exists j\in J\;u=s(i,j)\bigr)\vdash{}\\[-2pt]
    &(\exists u\in U\;F(u))\leftrightarrow
      (\exists i\in I\,\exists j\in J\;F(s(i,j)))
  \end{aligned}
}{TermImageWitnessSubstitution}{%
  \DeltaRow{Mengen}{U\dsep I\dsep J\dsep u\dsep i\dsep j}
  \DeltaRow{Mengenterm}{s(i,j)}
  \DeltaRow{Prädikat}{F(u)}
}

Zum Beweis wird zunächst die Vertauschung zweier unbeschränkter
Existenzquantoren in beiden Richtungen festgehalten. Dabei ist
\(H(a,b)\) ein beliebiges Prädikat mit frischen Bindungsvariablen
\(a,b\). In den anschließenden Zeilen steht
\(G(i,j)\) als bloße Abkürzung für \(i\in I\land j\in J\).
Beschränkte Quantoren werden durch ihre Mitgliedschaftsbedingungen
entfaltet; die Konjunktionen dürfen dabei nach dem Kommutativ- und
Assoziativgesetz umgeordnet werden.

\begin{tabproofsplitwide}
  \proofpartwide{Vertauschung der Existenzquantoren}
  \proofstepwidestar[1]{\exists a\,\exists b\;H(a,b)}{\rA}
  \proofstepwidestar[1]{\exists b\,\exists a\;H(a,b)}{%
    \FormulaRefAuto{\exists x\exists yF(x,y)\vdash
      \exists y\exists x F(x,y)}[thm]{1}}
  \proofstepwidestar[]{(\exists a\,\exists b\;H(a,b))\rightarrow
    (\exists b\,\exists a\;H(a,b))}{\rRI{1,2}}
  \proofstepwidestar[4]{\exists b\,\exists a\;H(a,b)}{\rA}
  \proofstepwidestar[4]{\exists a\,\exists b\;H(a,b)}{%
    \FormulaRefAuto{\exists x\exists yF(x,y)\vdash
      \exists y\exists x F(x,y)}[thm]{4}}
  \proofstepwidestar[]{(\exists b\,\exists a\;H(a,b))\rightarrow
    (\exists a\,\exists b\;H(a,b))}{\rRI{4,5}}
  \proofstepwidestar[]{(\exists a\,\exists b\;H(a,b))\leftrightarrow
    (\exists b\,\exists a\;H(a,b))}{\rLRI{3,6}}
  \closeproofpartwide

  \proofpartwide{Auswertung der Zwischenmenge}
  \setcounter{proofstepnr}{7}
  \proofstepwidestar[8]{\forall u\,\bigl(u\in U\leftrightarrow
    \exists i\in I\,\exists j\in J\;u=s(i,j)\bigr)}{\rA}
  \proofstepwidestar[8]{\begin{aligned}[t]
    &(\exists u\in U\;F(u))\leftrightarrow{}\\[-2pt]
    &\quad\exists u\,\bigl((\exists i\in I\,\exists j\in J\;
      u=s(i,j))\land F(u)\bigr)
  \end{aligned}}{\rLRS{\rUE{8},\FormulaRefAuto{P\leftrightarrow P}[thm]}}
  \proofstepwidestar[8]{\begin{aligned}[t]
    &(\exists u\in U\;F(u))\leftrightarrow{}\\[-2pt]
    &\quad\exists u\,\exists i\,\exists j\;
      (G(i,j)\land(u=s(i,j)\land F(u)))
  \end{aligned}}{%
    \rLRS{\FormulaRefAuto{\exists x(P \land Q(x))\eqvdash
      P \land \exists x(Q(x))}[thm],
      \FormulaRefAuto{P\land Q\eqvdash Q\land P}[thm],
      \FormulaRefAuto{P \land (Q \land R)\eqvdash
        (P \land Q) \land R}[thm],9}}
  \proofstepwidestar[8]{\begin{aligned}[t]
    &(\exists u\in U\;F(u))\leftrightarrow{}\\[-2pt]
    &\quad\exists i\,\exists j\,\exists u\;
      (G(i,j)\land(u=s(i,j)\land F(u)))
  \end{aligned}}{\rLRS{7,10}}
  \proofstepwidestar[]{F(s(i,j))\leftrightarrow
    \exists u\,(u=s(i,j)\land F(u))}{%
    \FormulaRefAuto{Fa \eqvdash \exists x(x=a \land Fx)}[thm]}
  \proofstepwidestar[]{\begin{aligned}[t]
    &\exists u\,(G(i,j)\land(u=s(i,j)\land F(u)))\\[-2pt]
    &\quad\leftrightarrow(G(i,j)\land F(s(i,j)))
  \end{aligned}}{%
    \rLRS{12,\FormulaRefAuto{\exists x(P \land Q(x))\eqvdash
      P \land \exists x(Q(x))}[thm]}}
  \proofstepwidestar[8]{(\exists u\in U\;F(u))\leftrightarrow
    (\exists i\in I\,\exists j\in J\;F(s(i,j)))}{%
    \rLRS{13,\FormulaRefAuto{\exists x(P \land Q(x))\eqvdash
      P \land \exists x(Q(x))}[thm],11}}
  \closeproofpartwide
\end{tabproofsplitwide}

Zeile 11 verwendet Zeile 7 zweimal: zunächst für \(u,i\), dann für
\(u,j\) unter dem äußeren Quantor \(\exists i\). Zeile 10 verwendet
\FormulaRefAuto{\exists x(P \land Q(x))\eqvdash
  P \land \exists x(Q(x))}[thm] jeweils für die von der gebundenen
Variablen unabhängigen Konjunkte. Damit ist auch die Umformung der
beschränkten Quantoren begründet; eine eigene Vertauschungsregel für
beschränkte Quantoren wird nicht vorausgesetzt.

Die übliche Bildformel ist eine Instanz dieses Satzes:
\[
  \{t(u)\mid u\in\{s(j)\mid j\in J\}\}
    =\{t(s(j))\mid j\in J\}.
\]
Hier wird \(U=\{s(j)\mid j\in J\}\) gewählt; der erste Index läuft
über \(J\), der zweite über die Einermenge \(\{\varnothing\}\),
und der erzeugende Term hängt nur vom ersten Index ab. Mit
\(F(u)\) als \(w=t(u)\) liefert
\FormulaRefAuto{TermImageWitnessSubstitution}[thm] die Äquivalenz der
beiden Elementkriterien aus
\FormulaRefAuto{B27TermImageMembership}[thm]. Der zusätzliche
Einermengenindex entfällt nach
\FormulaRefAuto{x\in\{a\}\eqvdash x=a}[thm] und
\FormulaRefAuto{a\in\{a\}}[thm]. Extensionalität ergibt die
angezeigte Gleichheit. Auch bei leerem \(J\) ist kein Zeuge zu wählen.

\medskip
Die beiden folgenden Sätze werten eine innere Termbildmenge unmittelbar
im ersten oder zweiten Argument eines weiteren Terms aus.
Dabei sind \(f(a,b)\) und \(g(a,b)\) Mengenterme mit festen Parametern.
Die verschiedenen Bindungsvariablen \(x,y,z,u,v,w\) kommen weder in
\(X,Y,Z\) noch in diesen Parametern frei vor; alle Einsetzungen sind
kollisionsfrei. Die Indexmengen dürfen leer sein.
Die Mehrfachindex- und Operatorschreibweise ist bei
\FormulaRefAuto{B27TermImageDef}[def] eingeführt.
Die zugehörigen Paarindizes liegen in den entsprechenden kartesischen
Produkten. Nach
\FormulaRefAuto{(a,b)\in A\times B\eqvdash a\in A\land b\in B}[thm],
\FormulaRefAuto{((a,b),c)\in (A\times B)\times C\ \eqvdash\ a\in A\land b\in B\land c\in C}[thm]
und den gebündelten Existenzregeln ist
\FormulaRefAuto{B27TermImageMembership}[thm] dann mit den angezeigten
einzelnen Indexvariablen lesbar. Die Termbildmengen existieren nach
\FormulaRefAuto{TermImageSetUniqueExists}[thm].

\FormulaThmDeltaK[Auswertung einer Termbildmenge im zweiten Argument]{%
  \begin{aligned}[t]
    &\{f(x,v)\mid x\in X,\,
      v\in\{g(y,z)\mid y\in Y,z\in Z\}\}\\[-2pt]
    &\qquad=\{f(x,g(y,z))\mid x\in X,y\in Y,z\in Z\}
  \end{aligned}
}{TermImageCompositionRight}{%
  \DeltaRow{Mengen}{X\dsep Y\dsep Z\dsep x\dsep y\dsep z\dsep v\dsep w}
  \DeltaRow{Mengenterme}{f(a,b)\dsep g(a,b)}
}

Im Beweis sei \(U=\{g(y,z)\mid y\in Y,z\in Z\}\);
\(M\) und \(N\) bezeichnen die linke und die rechte Menge der Aussage.
\begin{tabproofwide}
  \proofstepwidestar[]{\begin{aligned}[t]
    &\forall v\,(v\in U\leftrightarrow{}\\[-2pt]
    &\quad\exists y\in Y\,\exists z\in Z\;v=g(y,z))
  \end{aligned}}{\rUI{\FormulaRefAuto{B27TermImageMembership}[thm]}}
  \proofstepwidestar[]{\begin{aligned}[t]
    &w\in M\leftrightarrow{}\\[-2pt]
    &\quad\exists x\in X\,\exists v\in U\;w=f(x,v)
  \end{aligned}}{\FormulaRefAuto{B27TermImageMembership}[thm]}
  \proofstepwidestar[]{\begin{aligned}[t]
    &w\in M\leftrightarrow{}\\[-2pt]
    &\quad\exists x\in X\,\exists y\in Y\,\exists z\in Z\;
      w=f(x,g(y,z))
  \end{aligned}}{\rLRS{%
    \FormulaRefAuto{TermImageWitnessSubstitution}[thm]{1},2}}
  \proofstepwidestar[]{\begin{aligned}[t]
    &w\in N\leftrightarrow{}\\[-2pt]
    &\quad\exists x\in X\,\exists y\in Y\,\exists z\in Z\;
      w=f(x,g(y,z))
  \end{aligned}}{\FormulaRefAuto{B27TermImageMembership}[thm]}
  \proofstepwidestar[]{M=N}{%
    \FormulaRefAuto{\forall x\,(x\in A\leftrightarrow x\in B)
      \eqvdash A=B}[thm]{\rUI{\rLRS{4,3}}}}
\end{tabproofwide}

\Needspace{32\baselineskip}
\FormulaThmDeltaK[Auswertung einer Termbildmenge im ersten Argument]{%
  \begin{aligned}[t]
    &\{f(u,z)\mid u\in\{g(x,y)\mid x\in X,y\in Y\},\,
      z\in Z\}\\[-2pt]
    &\qquad=\{f(g(x,y),z)\mid x\in X,y\in Y,z\in Z\}
  \end{aligned}
}{TermImageCompositionLeft}{%
  \DeltaRow{Mengen}{X\dsep Y\dsep Z\dsep x\dsep y\dsep z\dsep u\dsep w}
  \DeltaRow{Mengenterme}{f(a,b)\dsep g(a,b)}
}

Im Beweis sei \(U=\{g(x,y)\mid x\in X,y\in Y\}\);
\(M\) und \(N\) bezeichnen wieder die linke und die rechte Menge.
\begin{tabproofwide}
  \proofstepwidestar[]{\begin{aligned}[t]
    &\forall u\,(u\in U\leftrightarrow{}\\[-2pt]
    &\quad\exists x\in X\,\exists y\in Y\;u=g(x,y))
  \end{aligned}}{\rUI{\FormulaRefAuto{B27TermImageMembership}[thm]}}
  \proofstepwidestar[]{\begin{aligned}[t]
    &w\in M\leftrightarrow{}\\[-2pt]
    &\quad\exists u\in U\,\exists z\in Z\;w=f(u,z)
  \end{aligned}}{\FormulaRefAuto{B27TermImageMembership}[thm]}
  \proofstepwidestar[]{\begin{aligned}[t]
    &w\in M\leftrightarrow{}\\[-2pt]
    &\quad\exists x\in X\,\exists y\in Y\,\exists z\in Z\;
      w=f(g(x,y),z)
  \end{aligned}}{\rLRS{%
    \FormulaRefAuto{TermImageWitnessSubstitution}[thm]{1},2}}
  \proofstepwidestar[]{\begin{aligned}[t]
    &w\in N\leftrightarrow{}\\[-2pt]
    &\quad\exists x\in X\,\exists y\in Y\,\exists z\in Z\;
      w=f(g(x,y),z)
  \end{aligned}}{\FormulaRefAuto{B27TermImageMembership}[thm]}
  \proofstepwidestar[]{M=N}{%
    \FormulaRefAuto{\forall x\,(x\in A\leftrightarrow x\in B)
      \eqvdash A=B}[thm]{\rUI{\rLRS{4,3}}}}
\end{tabproofwide}

Die folgenden Fassungen nehmen die Darstellung der Zwischenmenge als
Gleichheitsprämisse auf. Dabei sind \(f(a,b)\) und \(g(a,b)\) beliebige
Mengenterme mit festen Parametern. Die jeweils verwendeten
Bindungsvariablen \(x,y,z,u,v\) sind paarweise verschieden und kommen
weder in \(X,Y,Z,U,V\) noch in diesen festen Parametern frei vor.
Alle Einsetzungen sind kollisionsfrei. Die Indexmengen dürfen leer sein.

% Die neuen Deklarationen verwenden standalone ihre nativen Theorem-Anker,
% damit die bestehenden externen section*-Verweisziele erhalten bleiben.
\begingroup
\ifSubfilesClassLoaded{\let\phantomsection\relax}{}
\Needspace{24\baselineskip}
\FormulaThmDeltaK[Auswertung einer dargestellten Zwischenmenge im ersten Argument]{%
  \begin{aligned}[t]
    &U=\{g(x,y)\mid x\in X,y\in Y\}\vdash{}\\[-2pt]
    &\{f(u,z)\mid u\in U,z\in Z\}\\[-2pt]
    &\qquad=\{f(g(x,y),z)\mid x\in X,y\in Y,z\in Z\}
  \end{aligned}
}{TermImageCompositionLeftFromEquality}{%
  \DeltaRow{Mengen}{X\dsep Y\dsep Z\dsep U\dsep x\dsep y\dsep z\dsep u}
  \DeltaRow{Mengenterme}{f(a,b)\dsep g(a,b)}
}
\endgroup
\begin{tabproofwide}
  \proofstepwidestar[1]{U=\{g(x,y)\mid x\in X,y\in Y\}}{\rA}
  \proofstepwidestar[1]{\begin{aligned}[t]
    &\{f(u,z)\mid u\in U,z\in Z\}\\[-2pt]
    &\qquad=\{f(g(x,y),z)\mid x\in X,y\in Y,z\in Z\}
  \end{aligned}}{%
    \rIE{\FormulaRefAuto{a=b\vdash b=a}[thm]{1},
    \FormulaRefAuto{TermImageCompositionLeft}[thm]}}
\end{tabproofwide}

\begingroup
\ifSubfilesClassLoaded{\let\phantomsection\relax}{}
\Needspace{24\baselineskip}
\FormulaThmDeltaK[Einsetzen einer dargestellten Zwischenmenge im zweiten Argument]{%
  \begin{aligned}[t]
    &V=\{g(y,z)\mid y\in Y,z\in Z\}\vdash{}\\[-2pt]
    &\{f(x,g(y,z))\mid x\in X,y\in Y,z\in Z\}\\[-2pt]
    &\qquad=\{f(x,v)\mid x\in X,v\in V\}
  \end{aligned}
}{TermImageCompositionRightFromEquality}{%
  \DeltaRow{Mengen}{X\dsep Y\dsep Z\dsep V\dsep x\dsep y\dsep z\dsep v}
  \DeltaRow{Mengenterme}{f(a,b)\dsep g(a,b)}
}
\endgroup
\begin{tabproofwide}
  \proofstepwidestar[1]{V=\{g(y,z)\mid y\in Y,z\in Z\}}{\rA}
  \proofstepwidestar[1]{\begin{aligned}[t]
    &\{f(x,g(y,z))\mid x\in X,y\in Y,z\in Z\}\\[-2pt]
    &\qquad=\{f(x,v)\mid x\in X,v\in V\}
  \end{aligned}}{%
    \rIE{\FormulaRefAuto{a=b\vdash b=a}[thm]{1},
    \FormulaRefAuto{a=b\vdash b=a}[thm]{
    \FormulaRefAuto{TermImageCompositionRight}[thm]}}}
\end{tabproofwide}

\clearpage
\section{Zielmengen binärer Termbilder}

Der folgende Satz gilt für beliebige Mengenterme \(t(x,y)\), auch für
Funktionswerte \(F(x,y)\). Die Wertebedingung besagt dann, dass die
betrachtete Abbildung auf \(U\times V\) Werte in \(Z\) annimmt.
Die Teilmengen \(X\) und \(Y\) dürfen leer sein.

\FormulaThmDeltaK[Zielmenge eines binären Termbildes auf Teilmengen]{%
  \begin{aligned}[t]
    &X\subseteq U\dsep Y\subseteq V\dsep
      \forall x\in U\,\forall y\in V\;t(x,y)\in Z\vdash{}\\[-2pt]
    &\{t(x,y)\mid x\in X,y\in Y\}\subseteq Z
  \end{aligned}
}{BinaryTermImageSubset}{%
  \DeltaRow{Mengen}{U\dsep V\dsep X\dsep Y\dsep Z\dsep x\dsep y\dsep w}
  \DeltaRow{Mengenterm}{t(x,y)}
}
\begin{tabproofwide}
  \proofstepwidestar[1]{X\subseteq U}{\rA}
  \proofstepwidestar[2]{Y\subseteq V}{\rA}
  \proofstepwidestar[3]{\forall x\in U\,\forall y\in V\;t(x,y)\in Z}{\rA}
  \proofstepwidestar[4]{x\in X}{\rA}
  \proofstepwidestar[5]{y\in Y}{\rA}
  \proofstepwidestar[1,4]{x\in U}{%
    \FormulaRefAuto{A\subseteq B,\,x\in A\vdash x\in B}[thm]{1,4}}
  \proofstepwidestar[2,5]{y\in V}{%
    \FormulaRefAuto{A\subseteq B,\,x\in A\vdash x\in B}[thm]{2,5}}
  \proofstepwidestar[1,2,3,4,5]{t(x,y)\in Z}{%
    \rRE{\rUE{\rRE{\rUE{3},6}},7}}
  \proofstepwidestar[1,2,3]{\forall x\in X\,\forall y\in Y\;t(x,y)\in Z}{%
    \rUI{\rRI{4,\rUI{\rRI{5,8}}}}}
  \proofstepwidestar[1,2,3]{\{t(x,y)\mid x\in X,y\in Y\}\subseteq Z}{%
    \FormulaRefAuto{B27TermImageSubset}[thm]{9}}
\end{tabproofwide}

In der letzten Zeile wird die Mehrfachindexschreibweise aus
\FormulaRefAuto{B27TermImageDef}[def] als Paarindex gelesen.
Für binäre Operatoren ergibt sich daraus die unmittelbar verwendbare Fassung:

\Needspace{28\baselineskip}
\FormulaThmDeltaK[Binäre Operatoren bilden Teilmengen in den Träger ab]{%
  \begin{aligned}[t]
    &\mathsf{BinOp}(A,\star)\dsep X\subseteq A\dsep Y\subseteq A\vdash{}\\[-2pt]
    &\{x\star y\mid x\in X,y\in Y\}\subseteq A
  \end{aligned}
}{BinaryOperationImageSubset}{%
  \DeltaRow{Mengen}{A\dsep X\dsep Y\dsep x\dsep y}
  \DeltaRow{Funktionen}{\star}
}
\begin{tabproofwide}
  \proofstepwidestar[1]{\mathsf{BinOp}(A,\star)}{\rA}
  \proofstepwidestar[2]{X\subseteq A\land Y\subseteq A}{\rA}
  \proofstepwidestar[3]{x\in A}{\rA}
  \proofstepwidestar[4]{y\in A}{\rA}
  \proofstepwidestar[1,3,4]{x\star y\in A}{%
    \FormulaRefAuto{BinaryOperationClosureAxiom}[ax]{1,3,4}}
  \proofstepwidestar[1]{\forall x\in A\,\forall y\in A\;x\star y\in A}{%
    \rUI{\rRI{3,\rUI{\rRI{4,5}}}}}
  \proofstepwidestar[1,2]{\{x\star y\mid x\in X,y\in Y\}\subseteq A}{%
    \FormulaRefAuto{BinaryTermImageSubset}[thm]{2,2,6}}
\end{tabproofwide}

\Needspace{22\baselineskip}
\FormulaThmDeltaK[Binäre Termbilder gehören zur Potenzmenge des Trägers]{%
  \begin{aligned}[t]
    &\mathsf{BinOp}(A,\star)\dsep X\subseteq A\dsep Y\subseteq A\vdash{}\\[-2pt]
    &\{x\star y\mid x\in X,y\in Y\}\in\powerset(A)
  \end{aligned}
}{BinaryOperationImagePowerset}{%
  \DeltaRow{Mengen}{A\dsep X\dsep Y\dsep x\dsep y}
  \DeltaRow{Funktionen}{\star}
}
\begin{tabproofwide}
  \proofstepwidestar[1]{\mathsf{BinOp}(A,\star)}{\rA}
  \proofstepwidestar[2]{X\subseteq A\land Y\subseteq A}{\rA}
  \proofstepwidestar[1,2]{\{x\star y\mid x\in X,y\in Y\}\subseteq A}{%
    \FormulaRefAuto{BinaryOperationImageSubset}[thm]{1,2,2}}
  \proofstepwidestar[1,2]{\{x\star y\mid x\in X,y\in Y\}\in\powerset(A)}{%
    \FormulaRefAuto{x \in \powerset(A)\;\eqvdash\; x \subseteq A}[thm]{3}}
\end{tabproofwide}

Nichtleerheit benötigt nur die Nichtleerheit der beiden Indexmengen.
Der folgende Satz gilt deshalb allgemeiner für jeden Mengenterm \(t(x,y)\);
für einen binären Operator wird \(t(x,y)=x\star y\) eingesetzt.

\Needspace{28\baselineskip}
\FormulaThmDeltaK[Nichtleerheit eines binären Termbildes]{%
  X\neq\varnothing\dsep Y\neq\varnothing\vdash
    \{t(x,y)\mid x\in X,y\in Y\}\neq\varnothing
}{BinaryTermImageNonempty}{%
  \DeltaRow{Mengen}{X\dsep Y\dsep x\dsep y}
  \DeltaRow{Mengenterm}{t(x,y)}
}
\begin{tabproofwide}
  \proofstepwidestar[1]{X\neq\varnothing}{\rA}
  \proofstepwidestar[2]{Y\neq\varnothing}{\rA}
  \proofstepwidestar[3]{x\in X}{\rA}
  \proofstepwidestar[4]{y\in Y}{\rA}
  \proofstepwidestar[3,4]{t(x,y)\in\{t(u,v)\mid u\in X,v\in Y\}}{%
    \FormulaRefAuto{B27TermImageMembership}[thm]{\rEIStar{3,4,\rII}}}
  \proofstepwidestar[3,4]{\{t(x,y)\mid x\in X,y\in Y\}\neq\varnothing}{%
    \FormulaRefAuto{a\in A\vdash A\neq\varnothing}[thm]{5}}
  \proofstepwidestar[1,2]{\{t(x,y)\mid x\in X,y\in Y\}\neq\varnothing}{\begin{gathered}
    \RuleRef{EE}{\exists E}(\FormulaRefAuto{A\neq\varnothing\eqvdash\exists x\,(x\in A)}[thm]{1},3,\\[-2pt]
    \RuleRef{EE}{\exists E}(\FormulaRefAuto{A\neq\varnothing\eqvdash\exists x\,(x\in A)}[thm]{2},4,6))
  \end{gathered}}
\end{tabproofwide}

Die beiden Ergebnisse lassen sich für nichtleere Teilmengen zusammenfassen.

\Needspace{30\baselineskip}
\FormulaThmDeltaK[Binäre Operatoren bilden nichtleere Teilmengen auf nichtleere Teilmengen ab]{%
  \begin{aligned}[t]
    &\mathsf{BinOp}(A,\star)\dsep X,Y\in\PowNE{A}\vdash{}\\[-2pt]
    &\{x\star y\mid x\in X,y\in Y\}\in\PowNE{A}
  \end{aligned}
}{BinaryOperationImageNonemptyPowerset}{%
  \DeltaRow{Mengen}{A\dsep X\dsep Y\dsep x\dsep y}
  \DeltaRow{Funktionen}{\star}
}
\begin{tabproofwide}
  \proofstepwidestar[1]{\mathsf{BinOp}(A,\star)}{\rA}
  \proofstepwidestar[2]{X\in\PowNE{A}\land Y\in\PowNE{A}}{\rA}
  \proofstepwidestar[1,2]{\{x\star y\mid x\in X,y\in Y\}\in\powerset(A)}{\begin{gathered}
    \FormulaRefAuto{BinaryOperationImagePowerset}[thm](1,\\[-2pt]
    \FormulaRefAuto{NonemptyPowersetMembership}[thm]{2},\\[-2pt]
    \FormulaRefAuto{NonemptyPowersetMembership}[thm]{2})
  \end{gathered}}
  \proofstepwidestar[2]{\{x\star y\mid x\in X,y\in Y\}\neq\varnothing}{\begin{gathered}
    \FormulaRefAuto{BinaryTermImageNonempty}[thm](\\[-2pt]
    \FormulaRefAuto{NonemptyPowersetMembership}[thm]{2},\\[-2pt]
    \FormulaRefAuto{NonemptyPowersetMembership}[thm]{2})
  \end{gathered}}
  \proofstepwidestar[1,2]{\{x\star y\mid x\in X,y\in Y\}\in\PowNE{A}}{\begin{gathered}
    \FormulaRefAuto{NonemptyPowersetMembership}[thm](\\[-2pt]
    \RuleRef{AI}{\land I}(\\[-2pt]
    \FormulaRefAuto{x \in \powerset(A)\;\eqvdash\; x \subseteq A}[thm]{3},4))
  \end{gathered}}
\end{tabproofwide}
